\subsection{Relations with the support}

Let M be a module over a ring A. Recall that the set of prime ideals p of A such that \(M_{p} \neq 0\) is called the support of M and is denoted by Supp(M) (Chapter II, §4, no. 4, Definition 5).

PROPOSITION 7. Let A be a ring and M an A-module.

\begin{enumerate}
\def\labelenumi{(\roman{enumi})}
\item
  Every prime ideal \(\mathfrak{p}\) of A containing an element of \(\operatorname{Ass}(\mathbf{M})\) belongs to \(\operatorname{Supp}(\mathbf{M})\) .
\item
  Conversely, if A is Noetherian, every ideal \(\mathfrak{p} \in \operatorname{Supp}(\mathbf{M})\) contains an element of \(\operatorname{Ass}(\mathbf{M})\) .
\end{enumerate}

If \(\mathfrak{p}\) contains an element \(\mathfrak{q}\) of \(\operatorname{Ass}(\mathbf{M})\) , then \(\mathfrak{q} \cap (\mathbf{A} - \mathfrak{p}) = \varnothing\) and hence, if we write \(\mathbf{S} = \mathbf{A} - \mathfrak{p}\) , \(\mathbf{S}^{-1}\mathfrak{p}\) is a prime ideal associated with \(\mathbf{S}^{-1}\mathbf{M} = \mathbf{M}_{\mathfrak{p}}\) (no. 2, Proposition 5) and a fortiori \(\mathbf{M}_{\mathfrak{p}} \neq 0\) , hence \(\mathfrak{p} \in \operatorname{Supp}(\mathbf{M})\) . Conversely, if \(\mathbf{A}\) is Noetherian, so is \(\mathbf{A}_{\mathfrak{p}}\) (Chapter II, § 2, no. 4, Corollary 2 to Proposition 10). If \(\mathbf{M}_{\mathfrak{p}} \neq 0\) , then \(\operatorname{Ass}_{\mathbf{A}_{\mathfrak{p}}}(\mathbf{M}_{\mathfrak{p}}) \neq 0\) (no. 1, Corollary 1 to Proposition 2) and hence there exists \(\mathfrak{q} \in \operatorname{Ass}_{\mathbf{A}}(\mathbf{M})\) such that \(\mathfrak{q} \cap (\mathbf{A} - \mathfrak{p}) = \varnothing\) (no. 2, Corollary to Proposition 5).

COROLLARY 1. If M is a module over a Noetherian ring, then \(\operatorname{Ass}(\mathbf{M}) \subset \operatorname{Supp}(\mathbf{M})\) and these two sets have the same minimal elements.

COROLLARY 2. The nilradical of a Noetherian ring A is the intersection of the ideals \(\mathfrak{p} \in \operatorname{Ass}(\mathbf{A})\) .

We know that the nilradical of A is the intersection of the minimal elements of \(\operatorname{Spec}(A) = \operatorname{Supp}(A)\) (Chapter II, § 2, no. 6, Proposition 13).

